\section{Certificate formats and executable entry points}
\label{app:formats}

\subsection{Canonical rational polynomials}
A rational is encoded as \code{[numerator, denominator]}, with positive
 denominator and relatively prime integers. Booleans are not accepted as
integers, even though Python's type hierarchy normally allows that coercion.
A polynomial is a lexicographically sorted list of nonzero terms
\begin{lstlisting}
[ [exponent_vector, [numerator, denominator]], ... ]
\end{lstlisting}
with unique monomials. The empty list is zero. Exponent vectors have the
externally specified arity. The decoder rejects duplicate JSON keys,
floating-point values, unreduced rationals, explicit zero coefficients, and
unknown fields at mathematical boundaries.

The current implementation limits raw polynomial lists to 2,048 terms, encoded
total degree to 64, and rational numerator/denominator sizes to 4,096 bits.
Intermediate sparse arithmetic has a separate 20,000-term cap. A machine has
at most 16 state variables, 3 payload variables, 16 control nodes, and 64 edges;
a certificate basis has at most 128 entries per node. Default search limits
are smaller in several dimensions. These engineering limits are not part of
the mathematical completeness theorem.

These checks are not a hardened sandbox. Symbolic expansion, JSON nesting,
big-integer work, and adversarial expression sizes still require process-level
memory/time isolation before use on untrusted external inputs.

\subsection{Machine and positive certificate}
The source problem fixes its state and input dimensions, initial node and
state, all original edges and updates, and output goals at every control node.
Its semantics tag is exactly \code{rational-unrestricted-input}. The certificate
cannot change those choices.
\begin{lstlisting}
{
  "kind": "observable_closure",
  "basis": [ /* one polynomial list per node */ ],
  "goal_coordinates": [ /* every original goal */ ],
  "edge_coordinates": [ /* every original edge, in original order */ ]
}
\end{lstlisting}
Each edge row lists payload exponent vectors and rational coordinate vectors
in the source basis. The checker reconstructs the whole polynomial identity.
Omitting a row, edge, or goal does not weaken the task: it makes the certificate
invalid. Basis independence is deliberately not required for soundness.

\subsection{Execution certificate}
\begin{lstlisting}
{
  "kind": "execution_counterexample",
  "steps": [ {"edge": 0, "input": [[0,1]]}, ... ],
  "goal_index": 0
}
\end{lstlisting}
Replay starts from the original initial state. It checks edge applicability by
control node and evaluates all updates simultaneously. The indexed final goal
must be nonzero. There is no trusted final-state field, no user-supplied
residual value, and no need to trust the producer's invariant search.

\subsection{IPC problem and countermodel}
The problem has fields \code{version}, \code{kind}, \code{logic},
\code{context}, and \code{goal}, with \code{logic = "IPC"}. Formula nodes are
\code{["var", name]}, \code{["bot"]}, or a three-entry list beginning with
\code{"and"}, \code{"or"}, or \code{"imp"}.
\begin{lstlisting}
{
  "kind": "finite_kripke_countermodel",
  "worlds": 2,
  "root": 0,
  "relation": [[true,true],[false,true]],
  "valuation": {"P": [false,true]}
}
\end{lstlisting}
This particular model refutes excluded middle as an IPC formula. The checker
accepts at most 32 worlds; the recorded proposer uses at most 3. A changed logic
tag is rejected instead of being interpreted as a compatible proof system.

\subsection{Reproduce the experiment}
From the archive's root directory:
\begin{lstlisting}
python -m pip install -r requirements.txt
python -m pytest -q
python -S replay.py
python -m forge_delta.experiment --output reproduced-results --repeats 3
python emit_lean.py
\end{lstlisting}
The experiment refuses to overwrite a nonempty output directory. Its timings
will vary across machines; its deterministic problems and accepted certificate
contents should not. The full corpus is available directly to callers that
prefer the Python API:
\begin{lstlisting}
from forge_delta.cases import machine_cases
from forge_delta.search import solve
from forge_delta.checker import verify_positive

case = machine_cases()[0]
result = solve(case["problem"], case["limits"])
assert result["status"] == "proved"
assert verify_positive(case["problem"], result["certificate"])
\end{lstlisting}
Here ``proved'' is the prototype's machine-property status. It does not mean
that a Lean proof was constructed. For portable standalone replay use the
explicit checker rather than interpreting status strings as authority.

\subsection{Build the article and attempt the Lean gates}
\begin{lstlisting}
cd article
pdflatex -interaction=nonstopmode -halt-on-error forge_delta.tex
pdflatex -interaction=nonstopmode -halt-on-error forge_delta.tex
\end{lstlisting}
The second pass resolves cross-references and the table of contents. The
archive includes all article source sections. It does not redistribute fonts
or third-party repositories.

For Lean, use an existing project with the intended pinned Mathlib revision,
then run \code{lake env lean} on each supplied file using its absolute path.
The repository's recorded target was Lean \code{v4.34.0} with Mathlib
\code{1cf325a0cf67aca2b04d76b5380ff6a9e410aefa}. This package did not download,
compile, or certify that toolchain. The provided instructions do not silently
install dependencies or assert that generated files already pass.
